\SourceLeaf{153}Wahrscheinlichkeitsrechnung beantwortet werden soll. Nur kann ich kaum einräumen, das, was man auf diese Weise erhält, den \textls{wahrscheinlichsten} Wert zu nennen. Die Aufgabe, den wahrscheinlichsten Wert zu finden, ist tatsächlich mathematisch ganz bestimmt. Ihrer Natur nach setzt sie jedoch die Kenntnis des Fehlergesetzes voraus und führt \textls{nur} in dem einzigen Fall zu den arithmetischen Mitteln, in dem das Gesetz durch die Form $e^{-kkxx}$ ausgedrückt wird. Allgemein gesprochen ist die Wahrscheinlichkeit des wahrscheinlichsten Werts ebenfalls nur unendlich klein. Wenn $a$ der wahrscheinlichste Wert einer \textls{stetigen} Gesamtheit ist, bedeutet dies im Grunde nur, dass

\Needspace{7\baselineskip}
\begin{quote}
die Wahrscheinlichkeit, dass der wahre Wert zwischen $a-\omega$ und $a+\omega$ liegt, größer ist als die Wahrscheinlichkeit, dass der wahre Wert zwischen irgendeinem anderen Paar gleich weit auseinanderliegender Grenzen liegt, sofern $\omega$ unendlich klein ist.
\end{quote}

Genau betrachtet ist aber gerade deshalb ein solcher wahrscheinlichster Wert praktisch nur von geringem Interesse, weit weniger als derjenige Wert, bei dem der zu befürchtende Irrtum im Durchschnitt am wenigsten schädlich ist. Daher habe ich, abgesehen von anderen, freilich ebenso wichtigen oder noch viel wichtigeren Gründen, dieses zweite Prinzip vorgezogen, das mit dem ersten keinesfalls verwechselt werden darf.

. . . . . . . . . . . .

\Needspace{6\baselineskip}
\begin{center}\textsc{Gauss} an \textsc{Bessel}. Göttingen, 28. Februar 1839.\end{center}

. . . . . . . . . . . . . . . .

Ihren Aufsatz in den Astronomischen Nachrichten darüber, wie sich das Wahrscheinlichkeitsgesetz für Beobachtungsfehler aus zusammengesetzten Quellen der Formel $e^{-\frac{xx}{hh}}$ annähert, habe ich mit großem Interesse gelesen. Offen gesagt galt dieses Interesse weniger der Sache selbst als Ihrer Darstellung. Mit der Sache bin ich seit vielen Jahren vertraut, doch bin ich selbst nie dazu gekommen, die Entwicklung vollständig auszuführen.

